\documentclass[10pt,a4paper]{amsart}
\usepackage[margin=19mm]{geometry}
\usepackage{amsmath,amssymb,mathtools}
\usepackage[scr=rsfs,cal=euler]{mathalfa}
\usepackage{xfrac}
\usepackage[unicode,hidelinks,bookmarks=false]{hyperref}
\newcommand{\ie}{i.e.\ }
\newcommand{\cf}{cf.\ }
\newcommand{\resp}{respectively\ }
\newcommand{\Subsection}{\subsection}
\providecommand{\nobreakdash}{\nobreak\hbox{-}}
\setlength{\emergencystretch}{2em}
\multlinegap0pt
\usepackage{bm}
\usepackage{bbm}
\usepackage{dsfont}
\usepackage{xspace}
\usepackage{cancel}
\usepackage{mathtools}
\usepackage{amsthm}
\makeatletter
\@ifundefined{theorem}{\newtheorem{theorem}{Theorem}}{}
\makeatother
\newcommand {\N} {{\mathbb N}}
\newcommand {\Td} {{\mathbb T}^{d}}
\newcommand {\Z} {{\mathbb Z}}
\title[Orders of Oscillation Motivated by Sarnak's Conjecture--Part]{Orders of Oscillation Motivated by Sarnak's Conjecture--Part II}
\author{Yunping Jiang}
\date{Source: arXiv:2201.08800v2 (CC BY 4.0)}
\begin{document}
\maketitle

\noindent\emph{Source and licence.} Orders of Oscillation Motivated by Sarnak's Conjecture--Part II. Version arXiv:2201.08800v2 (CC BY 4.0), distributed under the Creative Commons Attribution 4.0 International licence, as recorded in the arXiv licence metadata for this exact version and shown on the abstract page https://arxiv.org/abs/2201.08800v2. Full rights notice: licenses/arxiv-2201.08800-CC-BY-4.0.txt.\par\smallskip

\noindent\textit{Editorial scope.} Counted unit: the Theorem whose source label is \texttt{main} in the article (source0.tex lines 347--349, source label \texttt{main}), delivered together with the article's own complete original proof of it (source0.tex lines 361--471, one \texttt{proof} environment of 111 lines). No mathematics is written, repaired or re-derived: statement and proof are the article's own text line for line. The proof is detached from the statement in the source: the article states the result at lines 347--349 and prints its proof at lines 361--471, and the proof environment's own header names the statement, so the association is the article's own. Required context retained uncounted: cn (lines 82--84); psp1 (lines 333--342). Every definition, display and label of those lines is retained unchanged. Omitted from this package and not redistributed: 1--81, 85--332, 343--346, 350--360, 472--966. Nothing inside a retained excerpt is omitted or altered: every retained body is delivered exactly as the article's lines are written, so no omission receipt of any kind accompanies this package. Citations: the retained excerpts carry no citation command at all, so no bibliography entry of the article is transcribed and no reference is redistributed. Reference adaptation: none was needed and none was made; every reference inside the delivered unit resolves inside it, so no unresolved reference or citation is produced. Printed slips kept literal: no printed word, symbol, hypothesis or number of the counted statement or of its proof was altered. Layout and class: the article's own class is \texttt{amsart} with packages including amsfonts, amssymb, amsmath, amsthm, enumerate, xy, graphicx, appendix, url, amscd, texdraw; the wrapper is the lane's pinned \texttt{amsart} editorial wrapper at 10pt on A4, which restates the theorem-like environments the retained text needs and adds the article's own macro definitions that the retained text needs (\texttt{\textbackslash N}, \texttt{\textbackslash Td}, \texttt{\textbackslash Z}), with extra wrapper packages (bm, bbm, dsfont, xspace, cancel, mathtools). The delivered numbering is the unit's own. Assets: no retained span contains a figure environment, a graphics-inclusion command, a PDF-page inclusion, a table or a TikZ picture; no image file is copied into this package, the package declares no asset dependency and no image file is redistributed with it. Licence: version arXiv:2201.08800v2 of this article is distributed under the Creative Commons Attribution 4.0 International licence, as recorded in the arXiv licence metadata for this exact version and shown on the abstract page https://arxiv.org/abs/2201.08800v2. A full rights notice is delivered at licenses/arxiv-2201.08800-CC-BY-4.0.txt. Characters: every retained excerpt is pure ASCII, so the delivered engine, XeLaTeX with its default Latin Modern text font, reports no missing character.
 \begin{equation}~\label{cn}
\frac{1}{N} \sum_{n=1}^{N} |c_{n}|^{\lambda} \leq C, \quad \forall\; N\geq 1.
\end{equation}  

\begin{equation}~\label{psp1}
f ({\bf x})= \left(
\begin{array}{c}
x_{1}+ a\\
x_{2} +h_{2} (x_{1})\\
x_{3} +b_{32}x_{2} +h_{3}(x_{1})\\
\vdots\\
x_{d}+b_{d(d-1)}x_{d-1} \cdots +b_{d2}x_{2}+h_{d} (x_{1})
\end{array}\right)
\end{equation}

\begin{theorem}~\label{main}
Suppose $d\geq 2$ and $k\geq 1$. Any oscillating sequence ${\bf c}$ of order $m=d+k-1$ and any simple polynomial skew products $f ({\bf x})$ of degree $k$ in (\ref{psp1}) are linearly disjoint.
\end{theorem}

\begin{proof}[Proof of Theorem~\ref{main}.]
Let ${\bf k}\in \Z^{d}$ with ${\bf k}^{t}=(k_{1}, \cdots,k_{d})$, define
$$
e({\bf k}\cdot {\bf x}) = e^{2\pi i (k_{1}x_{1}+\cdots + k_{d}x_{d})}.
$$
From the Stone-Weierstrass theorem, 
the set 
$S=\Big\{ e({\bf k}\cdot {\bf x})\Big\}_{{\bf k}\in \Z^{d}}$ 
forms a dense subset in $C(\Td)$.
A trigonometric polynomial $p$ is a linear combination of elements in $S$. We can write $p$ as 
$$
p({\bf x}) = \sum_{m_{1}\leq k_{1}\leq s_{1}}\cdots \sum_{m_{d}\leq k_{d}\leq s_{d}} a_{\bf k} e^{2\pi i (k_{1}x_{1}+\cdots +k_{d}x_{d})}.
$$
For any $\phi\in C(\Td)$, we have a sequence of trigonometric polynomials 
\begin{equation}~\label{stp}
p_{q}({\bf x}) = \sum_{m_{1q}\leq k_{1}\leq s_{1q}}\cdots \sum_{m_{dq}\leq k_{d}\leq s_{dq}} a_{{\bf k}, q} e^{2\pi i (k_{1}x_{1}+\cdots +k_{d}x_{d})}.
\end{equation}
such that $\|\phi-p_{q}\|\to 0$ as $q\to \infty$. The sequence $\{p_{q}\}_{q\in\N}$ is called the trigonometric approximation of $\phi$.
Consider
$$
S_{N}\phi ({\bf x}) = \frac{1}{N} \sum_{n=1}^{N} c_{n} \phi(f^{n} {\bf x}).
$$
 
For any $\epsilon>0$, we have an integer $r>0$ such that
$$
\| \phi-p_{r}\| <\frac{\epsilon}{2C^{\frac{1}{\lambda}}}.
$$
Then 
$$
S_{N}\phi ({\bf x}) =\Big( \frac{1}{N} \sum_{n=1}^{N} c_{n} \big( \phi(f^{n} {\bf x}) -p_{r} (f^{n} {\bf x})\big)\Big) +\Big( \frac{1}{N} \sum_{n=1}^{N} c_{n} p_{r} (f^{n} {\bf x})\Big) =I+II.
$$

For the estimation of $I$, we apply the H\"older inequality. Let $\lambda$ and $C$ be the numbers in (\ref{cn}) and $\lambda'>1$ be the dual number of $\lambda$, that is, 
$$
\frac{1}{\lambda} +\frac{1}{\lambda'}=1.
$$ 
Then we have that
$$
|I|\leq \Big( \frac{1}{N} \sum_{l=1}^{N} |c_{n}|^{\lambda}\Big)^{\frac{1}{\lambda}} 
 \Big( \frac{1}{N} \sum_{n=1}^{N} |\phi(f^{n} {\bf x})-p_{r}(f^{n} {\bf x})|^{\lambda'}\Big)^{\frac{1}{\lambda'}}
\leq C^{\frac{1}{\lambda}} \cdot \frac{\epsilon}{2C^{\frac{1}{\lambda}}} =\frac{\epsilon}{2}.
$$

For the estimation of $II$, let ${\bf x}_{n}=f^{n}{\bf x}$ for $n=1, 2,\cdots$.
Denote 
$$
{\bf x}_{n}^{t}=(x_{1}^{n},\cdots, x_{d}^{n}).
$$
Then 
$$
x_{1}^{n} = x_{1} + an
$$
is a polynomial of $n$ of degree less than or equal to $1$. Assume $x_{1}^{0}=x_{1}$. 

Suppose $h_{2} (x_{1}) = \sum_{i=0}^{k} r_{i} \cdot (x_{1})^{i}$. Then
$$
x_{2}^{n} = x_{2} +\sum_{l=0}^{n-1} h_{2}(x_{1}^{l})= x_{2}+\sum_{l=0}^{n-1} \sum_{i=0}^{k} r_{i} \cdot (x_{1}+la)^{i}
$$
$$
= \sum_{l=0}^{n-1} \sum_{i=0}^{k} s_{i} l^{i}
=\sum_{i=0}^{k} s_{i} \sum_{l=0}^{n-1} l^{i}= \sum_{i=0}^{k} s_{i} Q_{i}(n)
$$
where all $s_{i}$ are numbers only depending on $r_{i}$, $k$, $x_{1}$, and $a$ but not on $n$, and
$$
Q_{i}(n) = \sum_{l=0}^{n-1} l^{i}
$$ 
are polynomials of $n$ of degrees $i+1$ for $0\leq i\leq k$. 
Thus, $x_{2}^{n}$ is a polynomial of $n$ of degree less than or equal to $k+1$. Assume $x_{2}^{0}=x_{2}$.

Now
$$
x_{3}^{n} = x_{3} +b_{32} \Big( \sum_{l=0}^{n-1} x_{2}^{l} \Big) +\sum_{l=0}^{n-1} h_{3} (x_{1}^{l}).
$$
Similar to the calculation of $x_{2}^{n}$, we have that $\sum_{l=0}^{n-1} h_{3} (x_{1}^{l})$ is a polynomial of $n$ of degree less than or equal to $k+1$. Since $x_{2}^{n}$ is a polynomial of $n$ of degree less than or equal to $k+1$, $\sum_{l=0}^{n-1} x_{2}^{l}$ is a polynomial of $n$ of degree less than or equal to $k+2$. Thus, $x_{3}^{n}$ is a polynomial of $n$ of degree less than or equal to $k+2$. 
Inductively, we have that $x_{i}^{n}$ is a polynomial of $n$ of degree of $k+i-1$ for $4\leq i\leq d$. Furthermore,
$$
P_{{\bf k}} (n) = k_{1}x_{1}^{n}+\cdots +k_{d}x_{d}^{n}
$$
is a polynomial of $n$ of degree less than or equal to $d+k-1$ and 
$$
p_{r} (f^{n} {\bf x}) =  
\sum_{m_{1r} \leq k_{1}\leq s_{1r}} \cdots \sum_{ m_{dr} \leq k_{d}\leq s_{dr}} a_{{\bf k},r} e^{2\pi i P_{{\bf k}} (n)}.
$$

We can now continue to estimate 
$$
|II| =\Big|\frac{1}{N} \sum_{n=1}^{N} c_{n} \sum_{m_{1r} \leq k_{1}\leq s_{1r}} \cdots \sum_{m_{dr}\leq k_{d}\leq s_{dr}} a_{{\bf k},r} e^{2\pi i P_{{\bf k}} (n)}\Big| 
$$
$$
= \Big |\sum_{m_{1r}\leq k_{1}\leq s_{1r}}\cdots \sum_{m_{dr}\leq k_{d}\leq s_{dr}} a_{{\bf k},r} \frac{1}{N} \sum_{n=1}^{N} c_{n} e^{2\pi i P_{{\bf k}} (n)} \Big|.
$$
Let 
$$
L=\max \{ |m_{1r}|, \cdots, |m_{dr}|, |s_{1r}|, \cdots, |s_{dr}|, |a_{{\bf k},r}| \; ; \; m_{jr} \leq k_{j}\leq s_{jr}, 1\leq j\leq d\}.
$$
Since ${\bf c}$ is an oscillating sequence of order $d+k-1$, we can find an integer $M>r$ such that for $N>M$,
$$
\Big|\frac{1}{N} \sum_{n=1}^{N} c_{n} e^{2\pi i P_{{\bf k}} (n)} \Big|< \frac{\epsilon}{2L^{d}}, \quad \forall\; m_{1r} \leq k_{1}\leq s_{1r}, \cdots, m_{dr}\leq k_{n}\leq s_{dr}.
$$
This implies that 
$$
|II|< \epsilon/2.
$$

Combining the estimations of $I$ and $II$, we get that, for all $N>M$,
$$
|S_{N}\phi ({\bf x})|< \epsilon.
$$
This says that $\lim_{N\to \infty} S_{N}\phi({\bf x}) =0$. 
This completes the proof.
\end{proof}

\end{document}
